\section{Deformación positiva y dualidad reticular del registro generado}
\label{rec:positiva}

\subsection{Dirección de acoplamiento y flujo unimodular}

La extracción del registro, su lectura racional y los proyectores de incidencia se han construido en las secciones precedentes. Se conservan la base mil, el origen de fase, la frontera de acarreo y la coordenatización de las doce posiciones. En esa coordenatización,
\[
 K=(234,543,140,729,659,824,621,058,914,794,146,601),
 \qquad u=P_3P_{11}K,
\]
\[
 P_{11}=I-\frac1{12}\mathbf1\mathbf1^{\mathsf T},\qquad
 u\perp\mathbf1,\qquad
 \|u\|^2=\frac{6638585+2275584\sqrt5}{20}>0.
\]
La definición finita de \(P_3\), el orden de sus clases y la evaluación exacta de esta norma figuran en el lema~\ref{lem:k-sector-no-nulo}. La realización que sigue utiliza ese objeto generado. Escribimos
\begin{equation}
 \rho=\frac{\pi_{\HMT}}{\varphi_{\HMT}^2}>0,\qquad
 h_i=(\log\rho)\frac{u_i}{\|u\|},\qquad
 r_i(t)=e^{th_i},\qquad g_t=\operatorname{diag}(r_i(t)).
 \label{rec:eq:flujo}
\end{equation}
Las coordenadas regionales se evalúan antes de esta composición. El parámetro \(t\in\RR\) indexa una deformación matemática; su definición no lo identifica con el reloj TPK ni con una profundidad de refinamiento.

\begin{proposition}[Compensación de escala]
\label{rec:prop:flujo}
El flujo satisface
\[
 g_tg_s=g_{t+s},\qquad g_0=I,\qquad g_t^{-1}=g_{-t},\qquad
 \det g_t=\prod_i r_i(t)=1.
\]
\end{proposition}
\begin{proof}
La ortogonalidad al modo uniforme da \(\sum_i h_i=0\). La multiplicación de exponenciales prueba la ley de grupo y la inversión. El determinante es \(\exp(t\sum_i h_i)=1\). Cada factor diagonal es positivo.
\end{proof}

La conservación es multiplicativa y global: cada radio puede variar. En coordenadas logarítmicas se conserva la suma nula. Esta distinción será la que permita interpretar geométricamente el proyector transversal de rango diez.

\subsection{Reciprocidad exacta del doble círculo}
\label{rec:doble-circulo}

Consideremos la coordenatización compleja y su inversa
\[
 z(s)=\frac{s-1}{s},\qquad s(z)=\frac1{1-z}.
\]
Para \(r>0\), \(z=re^{i\theta}\ne1\), definimos \(s(r,\theta)=(1-re^{i\theta})^{-1}\).

\begin{proposition}[Conjugación del espejo espectral]
\label{rec:prop:espejo}
Se cumplen
\begin{align}
 s(r^{-1},\theta)&=1-\overline{s(r,\theta)},\label{rec:eq:espejo}\\
 \Re s(r,\theta)-\frac12&=
 \frac{1-r^2}{2|1-re^{i\theta}|^2}.\label{rec:eq:lado}
\end{align}
El círculo unidad, salvo el punto excluido \(z=1\), corresponde a la recta \(\Re s=1/2\).
\end{proposition}
\begin{proof}
La inversión radial a ángulo fijo es \(z'=1/\overline z\). Por tanto,
\[
 \frac1{1-z'}=\frac{\overline z}{\overline z-1}
 =1-\frac1{1-\overline z}.
\]
Además, \((1-z)^{-1}=(1-\overline z)/|1-z|^2\). Restar \(1/2\) de su parte real produce \eqref{rec:eq:lado}.
\end{proof}

La relación con los radios de \eqref{rec:eq:flujo} es explícita: \(r_i(-t)=r_i(t)^{-1}\). La igualdad de productos conserva volumen, mientras la proposición identifica una involución entre coordenatizaciones complejas. La localización de ceros de una función requiere las propiedades de esa función; no es una consecuencia del producto unitario de radios.

Conviene distinguir tres realizaciones circulares del corpus. La coordenatización espectral lleva \(s=1/2+i\gamma\) a \(z_\gamma=(\gamma+i/2)/(\gamma-i/2)\). El círculo dual de la memoria entera lleva un entero \(n\) al carácter \(z^n\). La compresión nonádica lleva \(z\) a \((8+z)/9\), círculo de centro \(8/9\) y radio \(1/9\), tangente al círculo unidad en \(1\). La razón modal \(5/9\) pertenece, en cambio, a los autovalores \(9,5\) del bloque APP \(\left(\begin{smallmatrix}7&2\\2&7\end{smallmatrix}\right)\). Sus funciones permanecen diferenciadas al componerlas.

\subsection{Estratos positroidales e invariantes equilibrados}

Sea \(X\) una matriz real de rango \(k\), con doce columnas, y sea \(\Delta_I(X)\) su menor maximal de columnas \(I\). La acción derecha satisface
\begin{equation}
 \Delta_I(Xg_t)=e^{t\sum_{i\in I}h_i}\Delta_I(X).
 \label{rec:eq:plucker}
\end{equation}
La positividad diagonal conserva todos los signos y todos los menores nulos. Así, para \(1\le k\le11\), una trayectoria iniciada en \(\operatorname{Gr}_{\ge0}(k,12)\) permanece en el mismo estrato positroidal. Se utiliza aquí la estratificación por menores como realización posterior; véase la construcción general de Postnikov~\cite{rec:postnikov}.

\begin{proposition}[Dualidad positiva de rangos complementarios]
Sea \(A_{\pm}=\operatorname{diag}(1,-1,1,-1,\ldots)\) y
\(\mathfrak D(V)=V^\perp A_{\pm}\), con la orientación positiva de sus menores complementarios. Entonces
\[
 \mathfrak D(Vg_t)=\mathfrak D(V)g_{-t}.
\]
\end{proposition}
\begin{proof}
Como \(g_t\) es diagonal y autoadjunto,
\((Vg_t)^\perp=V^\perp g_t^{-1}\). La matriz \(A_{\pm}\) conmuta con \(g_t\). Para el signo de los menores, el complemento ortogonal introduce, salvo un signo global, \((-1)^{\sum_{j\in J}j}\). La multiplicación de las columnas por \((-1)^{j-1}\) deja un signo independiente de \(J\); se elimina mediante una única orientación global. Quedan menores complementarios no negativos.
\end{proof}

Si \(\mathbf1_I+\mathbf1_J=\mathbf1_L+\mathbf1_M\) y el denominador es no nulo, la cancelación coordenada a coordenada de los exponentes prueba
\begin{equation}
 \frac{\Delta_I(Xg_t)\Delta_J(Xg_t)}{\Delta_L(Xg_t)\Delta_M(Xg_t)}
 =\frac{\Delta_I(X)\Delta_J(X)}{\Delta_L(X)\Delta_M(X)}.
 \label{rec:eq:cocientes}
\end{equation}
En rango dos se recuperan los cocientes de cuatro índices y sus relaciones de intercambio. En contraste, el producto de menores complementarios de las dos trayectorias emparejadas recibe dos veces el mismo factor y puede variar. El marco circulante \(O_K\), invertible para la recuperación, tampoco adquiere por esta acción la propiedad de ser totalmente no negativo: sus menores conservan los signos establecidos en su propia realización.

\subsection{Covolumen, red dual y simetría ternaria}

Sea \(\Lambda\) la red de Leech marcada construida anteriormente, autodual y de covolumen uno en la métrica de sus doce planos ortogonales de tipo \(A_2\). La isometría marcada
\[
 \sigma=\bigoplus_{i=1}^{12}P_i\mathsf C P_i^{-1},\qquad
 \mathsf C^2+\mathsf C+I=0,
\]
preserva el pegado y \(\Lambda\), tiene orden tres y carece de vectores fijos no nulos. Los \(P_i\) son marcos de los planos, distintos de los proyectores \(P_3,P_{11}\). Definimos
\begin{equation}
 G_t=\bigoplus_{i=1}^{12}r_i(t)I_{A_{2,i}},\qquad
 \Lambda_t=G_t\Lambda.
 \label{rec:eq:red}
\end{equation}

\begin{theorem}[Deformación reticular recíproca]
\label{rec:thm:red}
Para todo \(t\in\RR\),
\[
 \operatorname{covol}(\Lambda_t)=1,\qquad
 \Lambda_t^*=\Lambda_{-t},\qquad \sigma\Lambda_t=\Lambda_t.
\]
Toda isometría de \(\Lambda\) que conmuta con \(G_t\) preserva también \(\Lambda_t\).
\end{theorem}
\begin{proof}
Cada radio actúa sobre dos coordenadas; así,
\(\det G_t=\prod_i r_i(t)^2=1\). Para un operador invertible \(G\), un vector \(y\) pertenece a \((G\Lambda)^*\) exactamente cuando \(G^*y\in\Lambda^*\). Por ello
\((G\Lambda)^*=G^{-*}\Lambda^*\). Al aplicar autodualidad, autoadjunción e inversión se obtiene \(\Lambda_t^*=G_{-t}\Lambda\). Finalmente, \(G_t\) es escalar en cada plano y conmuta con \(\sigma\), de donde
\(\sigma G_t\Lambda=G_t\sigma\Lambda=G_t\Lambda\). El mismo cálculo sirve para las demás isometrías conmutantes. El operador \(\sigma\) no cambia, por lo que conserva su orden y su ausencia de vectores fijos.
\end{proof}

El punto \(t=0\) recupera la red original con su integralidad, paridad, autodualidad y realización excepcional. La colección completa es una colección de redes euclídeas unimodulares en el sentido de covolumen uno; fuera del punto original no se presupone integralidad ni una realización de Moonshine. El teorema identifica exactamente la simetría ternaria que sí persiste y la dualidad que relaciona los dos sentidos de deformación.

\begin{figure}[H]
\centering
\begin{tikzpicture}[node distance=3mm,>=Latex,every node/.style={font=\fontsize{8.5}{10}\selectfont,align=center},box/.style={draw=ink,rounded corners=2pt,text width=0.16\linewidth,inner sep=3pt,minimum height=25mm}]
\node[box] (a) {APP--TRIT--TPK: estado enriquecido y continuo conjunto};
\node[box,right=of a] (b) {Registro \(K\), representaciones regionales y archivo bilateral};
\node[box,right=of b] (c) {Dirección recuperable \(u_K\) y flujo positivo de determinante uno};
\node[box,right=of c] (d) {Radios conjugados, redes duales y coordenatizaciones hiperbólicas};
\node[box,right=of d] (e) {Reciprocidad espectral y recuperación orientada de la acción};
\draw[->] (a)--(b);\draw[->](b)--(c);\draw[->](c)--(d);\draw[->](d)--(e);
\end{tikzpicture}
\caption{Dependencias de la ampliación. Las flechas representan composiciones especificadas en las demostraciones; la recuperación conserva las coordenatizaciones y la orientación.}
\label{fig:rec-causal}
\end{figure}
